\documentclass[11pt,a4paper]{article}
\usepackage[spanish,es-noquoting]{babel}
\usepackage[top=2.2cm, bottom=2.2cm, left=2.5cm, right=2.5cm]{geometry}
\usepackage{fontspec}
\usepackage{amsmath,amssymb}
\usepackage{enumitem}
\usepackage{microtype}

\setmainfont{Calibri}
\linespread{1.15}

\begin{document}
\pagestyle{empty}

\begin{center}
    {\LARGE \textbf{Ejercicios de Recta Tangente y Normal}}\\[4pt]
    {\large \textbf{Matemáticas 1.º de Bachillerato}}
\end{center}

\vspace{0.4cm}
\noindent\textbf{1.} Dada la función polinómica $f(x) = x^3 - 4x^2 + x + 6$:
\begin{enumerate}[label=\textbf{\alph*)}, leftmargin=2em, itemsep=6pt]
    \item Halla la ecuación de la recta tangente a la gráfica de $f(x)$ en el punto de abscisa $x_0 = 3$.
    \item Halla la ecuación de la recta normal a la gráfica en dicho punto.
\end{enumerate}

\vspace{0.6cm}
\noindent\textbf{2.} Dada la función cuadrática $g(x) = x^2 - 6x + 5$:
\begin{enumerate}[label=\textbf{\alph*)}, leftmargin=2em, itemsep=6pt]
    \item Determina el punto de la gráfica en el que la recta tangente es paralela a la recta $r: y = 2x + 7$. Escribe la ecuación de dicha recta tangente.
    \item Encuentra el punto de la gráfica donde la recta tangente es horizontal (pendiente nula) y escribe su ecuación.
\end{enumerate}

\end{document}
